\documentclass[10pt,a4paper,sans]{moderncv}
\moderncvstyle{banking}
\moderncvcolor{blue}

% --- Шрифты и поддержка кириллицы ---
\usepackage{fontspec}
\usepackage{polyglossia}
\setdefaultlanguage{russian}
\setotherlanguage{english}

% Основные шрифты (можешь поменять на свои)
\setmainfont{Times New Roman}
\setsansfont{Arial}
\setmonofont{Courier New}

% Кириллица для polyglossia
\newfontfamily\cyrillicfont{Times New Roman}
\newfontfamily\cyrillicfontsf{Arial}
\newfontfamily\cyrillicfonttt{Courier New}

% --- Верстка ---
\usepackage[scale=0.885]{geometry}
\usepackage{enumitem}
\setlist[itemize]{leftmargin=*, noitemsep, topsep=0pt}

% --- Личные данные ---
\name{Нурсултан}{Куандык}
\extrainfo{n.kuandyk1995@gmail.com · +7 747 454 8448 · github.com/lunarnuts}
\vspace*{-1.0em}

\begin{document}
\makecvtitle
\setlength{\hintscolumnwidth}{2.2cm}
\renewcommand{\subsectionfont}{\normalsize\bfseries}
\setlength{\parskip}{0.2em}
\setlength{\parindent}{0pt}

% --- Профессиональное резюме ---
\section{Коротко обо мне}
Результативный \textbf{Senior Go Software Engineer} с опытом более 5 лет в проектировании и масштабировании облачных платформ, высокодоступных микросервисов, k8s операторов и event-driven систем в финтехе и корпоративном сегменте. Достиг сокращения уязвимостей на \textbf{89\%} (устранены все критические), увеличения покрытия тестами на \textbf{36\%} в кодовой базе объемом более 60k LOC, а также реализации API с нагрузкой 10k+ RPS и \textbf{uptime 99.9\%}. Опыт в согласовании архитектуры между командами, наставничестве и модернизации систем с глубокими знаниями Go, Kubernetes, gRPC, Kafka и PostgreSQL.

% --- Навыки ---
\section{Ключевые навыки}
\cvitem{Языки}{Go, JavaScript (Node.js, AngularJS), ABAP}
\cvitem{Архитектура}{Clean Architecture, Hexagonal, SOLID, DI; Event-driven и распределённые системы}
\cvitem{Облако/DevOps}{Kubernetes, Docker, AWS, GCP, ArgoCD, GitHub Actions, CI/CD, Operator SDK}
\cvitem{Базы данных}{PostgreSQL, MySQL, MongoDB, Redis; оптимизация запросов, транзакции, пакетная обработка}
\cvitem{Сообщения/API}{gRPC, Kafka, NATS, REST, OpenAPI 3.0}
\cvitem{Мониторинг}{Grafana, OpenTelemetry, Cribl, SignalFx, Splunk, Loki, Prometheus; инцидент-менеджмент, L3 поддержка}

% --- Опыт работы ---
\section{Опыт работы}

\cventry{Янв 2025 -- наст. время}{Senior Go Software Engineer}{EPAM}{София, Болгария}{}{
\begin{itemize}
  \item Принял ключевое участие в разработке операторного фреймворка для внутреннего использования.
  \item Улучшил наблюдаемость, упростив технический стек с использованием OpenTelemetry collector и Splunk(+SignalFx).
  \item Создал \textbf{Kubernetes операторов} для стандартизации деплойментов и rollout конфигураций, сократив ручные операции.
  \item Руководил командой из 3 инженеров при модернизации платформы на Node/Angular, сократил уязвимости с \textbf{228 (28 критических)} до \textbf{26 (0 критических)}.
  \item Увеличил покрытие тестами на \textbf{36\%} в кодовой базе из \textbf{60k+ LOC}.
  \item Спроектировал и реализовал масштабируемый \textbf{сервис многоканальных уведомлений} (in-app, email, SMS).
  \item \textbf{Организовал сессии проектирования между командами} для согласования API и архитектурных стандартов.
  \item Наставлял middle-разработчиков, проводил code review и обучал коллег Go и облачным практикам.
\end{itemize}}

\cventry{Янв 2024 -- Янв 2025}{Senior Go Software Engineer}{DigitalBizFactory}{Астана, Казахстан}{}{
\begin{itemize}
  \item Спроектировал и реализовал \textbf{core banking API} и микросервисы с worker pools и асинхронными планировщиками, обеспечил \textbf{10k RPS при uptime 99.9\%}.
  \item Внедрил \textbf{gRPC} и \textbf{Kafka} для декуплинга сервисов и потоковой передачи событий; использовал Redis для кэширования и rate-лимитов.
  \item Оптимизировал PostgreSQL запросы и индексы, снизил p95 latency на критических эндпоинтах на \textbf{60\%}.
  \item Согласовал SLA совместно с продукт-оунерами и QA.
\end{itemize}}

\cventry{Янв 2023 -- Янв 2024}{Middle Go Software Engineer}{Airba Group}{Алматы, Казахстан}{}{
\begin{itemize}
  \item Разработал безопасные финансовые API; автоматизировал аудит и сократил ручную подготовку отчетности с \textbf{40+ часов/мес} до \textbf{2 часов/мес}.
  \item Оптимизировал банковские интеграции с пакетной обработкой и реструктуризацией транзакций, снизил время ответа на \textbf{45\%}.
  \item Работал с кросс-функциональными командами для повышения безопасности данных.
\end{itemize}}

\cventry{Окт 2021 -- Янв 2023}{Junior $\rightarrow$ Middle Software Engineer}{EPAM}{Астана, Казахстан}{}{
\begin{itemize}
  \item Разработал backend API и интеграции с внешними системами для крупного нефтегазового клиента.
\end{itemize}}

% --- Образование ---
\section{Образование}
\cventry{2012 -- 2017}{Бакалавр Робототехники и Мехатроники}{Nazarbayev University}{Астана}{}{}

\end{document}